\section{Introducción}
\label{sec:introduccion-e4}

La Entrega 4 consolida el sistema distribuido de gestión de laboratorios como
un producto integrado. El trabajo parte de la infraestructura distribuida ya
disponible, pero su aportación no consiste en presentar nuevamente esa base: se
centra en delimitar los servicios, ordenar el código de negocio, ofrecer clientes
web y móvil y establecer mecanismos verificables de seguridad, pruebas,
integración y despliegue continuos y observabilidad. La arquitectura se trata,
por tanto, como un conjunto de decisiones que condicionan atributos de calidad
y no solamente como un diagrama de componentes. Esta perspectiva es coherente
con el papel que la literatura asigna a la arquitectura en la modificabilidad,
la seguridad, el rendimiento y la testabilidad de un sistema
\cite{bass2021softwarearchitecture}.

La solución incorpora una aplicación web en React y una aplicación Android con
Jetpack Compose, ambas consumidoras de un punto de entrada común. La seguridad
se aplica mediante autenticación JWT, sesiones de renovación y autorización por
roles; su centralización en el perímetro reduce exposiciones innecesarias, aunque
no sustituye los controles de cada servicio. Esta defensa en varios niveles es
relevante porque la comunicación distribuida amplía la superficie de ataque y
obliga a considerar conjuntamente identidad, infraestructura y desarrollo
\cite{berardi2022microservicesecurity}.

La verificación comprende pruebas unitarias, de integración, contratos,
interfaces y carga, coordinadas por el flujo de CI/CD. Prometheus, Grafana,
OpenTelemetry, Tempo, Loki y cAdvisor aportan métricas, trazas, registros y datos
de recursos para observar el comportamiento desplegado. Finalmente, la
evaluación adopta el modelo de calidad de producto ISO/IEC 25010:2023, que
proporciona características y subcaracterísticas para especificar y evaluar la
calidad de productos TIC \cite{iso25010_2023}. Los capítulos posteriores
distinguen las propiedades demostradas de las mediciones incompletas, evitando
convertir configuraciones o pruebas no ejecutadas en resultados experimentales.

\section{Arquitectura de la solución}
\label{sec:arquitectura-solucion}

\subsection{Límites de servicio y comunicación}
\label{subsec:limites-servicio}

El sistema adopta una descomposición por capacidades, formalizada en el
ADR-002\footnote{Evidencia interna: \path{../adr/ADR-002-microservicios.md}.}.
\texttt{auth-service} gestiona autenticación, emisión de tokens y sesiones de
renovación; \texttt{usuarios-service} administra perfiles institucionales;
\texttt{academico-laboratorios-service} mantiene los catálogos académicos y el
inventario de espacios y equipos; y \texttt{reservas-solicitudes-service}
controla solicitudes, disponibilidad, bloqueos y reservas. El
\texttt{api-gateway} completa la topología como punto de entrada para ruteo,
CORS y validación de autenticación, sin incorporar reglas del dominio.

Los clientes React y Android acceden a las API a través del Gateway. Las
operaciones que requieren información de otra capacidad utilizan contratos HTTP
explícitos entre servicios; no comparten repositorios ni entidades JPA. Cada
servicio conserva sus migraciones y su base lógica en CockroachDB. Esta
separación favorece la evolución y prueba independiente, pero introduce latencia,
errores parciales y compatibilidad de contratos como preocupaciones de diseño.
La literatura reporta precisamente que el despliegue y los patrones de
comunicación son tan relevantes como la delimitación inicial de los
microservicios \cite{aksakalli2021deployment}.

Docker proporciona unidades reproducibles de ejecución para los servicios y la
infraestructura de soporte. El flujo observable es: cliente, Gateway, servicio
propietario del caso de uso y persistencia. En paralelo, los servicios exportan
telemetría hacia la infraestructura de observabilidad. Los diagramas C4 de
contexto y contenedores existentes\footnote{Evidencia interna:
\path{../diagrams/c4-nivel1.png} y
\path{../diagrams/c4-nivel2.png}.} documentan estos límites; se referencian sin
duplicarlos dentro del nuevo informe.

La descomposición no elimina los compromisos propios de un sistema distribuido.
No existe una transacción ACID única que abarque las bases de todos los
servicios, de modo que los casos de uso interservicio deben gestionar fallos y
consistencia de forma explícita. Asimismo, el Gateway constituye un punto común
de acceso, pero los servicios siguen siendo responsables de validar las
operaciones que protegen. Estas consecuencias se consideran parte de la
arquitectura real y no defectos ocultos por la representación C4.

\subsection{Decisiones arquitectónicas registradas}
\label{subsec:adr-e4}

Los ADR preservan el contexto y las consecuencias de decisiones que el código
por sí solo no explica. El ADR-001\footnote{Evidencia interna:
\path{../adr/ADR-001-arquitectura.md}.} registra la reorganización por capas de
\texttt{academico-laboratorios-service}; el ADR-002 fija los límites de los
microservicios y la propiedad de sus datos. El ADR-005\footnote{Evidencia
interna: \path{../adr/ADR-005-patrones-gof.md}.} identifica los patrones que
participan en rutas reales y también documenta deuda arquitectónica. Por último,
el ADR-006\footnote{Evidencia interna:
\path{../adr/ADR-006-eleccion-movil.md}.} justifica Kotlin y Jetpack Compose,
MVVM, Room y las integraciones nativas como base del cliente móvil.

Los ADR-003 y ADR-004 no se presentan como resultados de esta entrega: sus
decisiones sobre distribución de datos constituyen antecedentes de la
infraestructura. Esta distinción evita atribuir a la Entrega 4 trabajo ya
documentado y mantiene el foco en la consolidación arquitectónica y de calidad.

\section{Arquitectura por capas y principios SOLID}
\label{sec:capas-solid}

\subsection{Organización del código}

Los servicios con reglas de negocio se organizan, con variaciones heredadas, en
cuatro zonas. La capa \emph{presentation} contiene controladores REST, DTO y la
traducción de excepciones a respuestas HTTP. La capa \emph{application}
orquesta servicios y casos de uso; decide el orden de las operaciones, pero
delega las reglas específicas. \emph{Domain} contiene modelos, estados,
estrategias, eventos y puertos que expresan las necesidades del negocio.
Finalmente, \emph{infrastructure} aporta adaptadores HTTP, seguridad,
observabilidad y persistencia, incluidas las entidades JPA y las
implementaciones Spring Data.

Esta dirección permite que un controlador dependa de un caso de uso y que el
caso de uso dependa de un puerto, mientras el adaptador concreto implementa ese
puerto desde el exterior. En consecuencia, una prueba puede reemplazar la base
de datos o un cliente remoto por un doble sin trasladar detalles de JPA al
escenario de negocio. El API Gateway es una excepción deliberada: sus paquetes
de configuración, rutas, seguridad y registro corresponden a infraestructura;
crear para él un dominio vacío no mejoraría la separación.

\subsection{Aplicación razonada de SOLID}

La responsabilidad única se refleja en controladores que adaptan HTTP,
servicios que coordinan casos de uso y repositorios que resuelven persistencia.
La inversión de dependencias se materializa cuando aplicación y dominio conocen
interfaces como \texttt{UsuarioRepository} o los puertos
\texttt{*RepositoryPort}, y los adaptadores de infraestructura dependen de esas
abstracciones. Strategy y State ofrecen puntos de extensión sin reescribir los
servicios coordinadores, lo que aproxima el principio abierto/cerrado. Las
interfaces son acotadas al comportamiento requerido por sus consumidores, y
las implementaciones de estados pueden sustituirse a través de sus contratos.

Estos elementos mejoran mantenibilidad y testabilidad, pero no demuestran un
cumplimiento perfecto de SOLID. En
\texttt{academico-laboratorios-service}, varios puertos del dominio importan
\texttt{org.springframework.data.domain.Page} y
\texttt{org.springframework.data.domain.Pageable}. El dominio queda así
acoplado al modelo de paginación de Spring Data, pese a que la persistencia se
implemente mediante adaptadores. Sustituir esos tipos por contratos propios y
mapearlos en infraestructura cerraría la brecha, pero exige una migración
coordinada de servicios y controladores. Se registra, por ello, como deuda
técnica y no se oculta bajo una afirmación absoluta de independencia.

\section{Patrones de diseño aplicados}
\label{sec:patrones-gof}

Los patrones siguientes se identificaron por sus participantes, consumidores y
comportamiento, no solamente por el nombre de una clase. Su uso responde a
problemas concretos del sistema y complementa la separación por capas.

\subsection{Repository y Factory Method}

\textbf{Repository.} El problema consiste en evitar que los casos de uso
conozcan JPA o Spring Data. En Auth, \texttt{UsuarioRepository} define el puerto
del dominio y \texttt{UsuarioPersistenceAdapter} lo implementa mediante
\texttt{UsuarioAuthRepository}; el mismo esquema puerto--adaptador aparece en
Usuarios, Académico y Reservas. El beneficio es sustituir la persistencia en
pruebas y concentrar el mapeo. Su aplicación es parcial en los puertos paginados
de Académico por la dependencia de Spring ya señalada.

\textbf{Factory Method.} La construcción del principal de seguridad no debe
quedar fijada en el servicio que carga usuarios. \texttt{UserDetailsFactory}
declara \texttt{crear(Usuario)} y \texttt{DefaultUserDetailsFactory} produce
\texttt{CustomUserDetails}; \texttt{CustomUserDetailsService} consume la
interfaz. La construcción queda encapsulada y puede sustituirse en pruebas sin
duplicar conocimiento del producto concreto.

\subsection{Strategy y State}

\textbf{Strategy.} La disponibilidad de un laboratorio es una política que
puede evolucionar sin reescribir la orquestación. La interfaz
\texttt{DisponibilidadStrategy} define la evaluación y
\texttt{DisponibilidadSinConflictosStrategy} comprueba fecha, laboratorio,
cruces y bloqueos. \texttt{StrategyConfig} selecciona la implementación y
\texttt{DisponibilidadServiceImpl} la utiliza. Aunque actualmente existe una
sola estrategia concreta, el punto de variación es real, puro y probado.

\textbf{State.} Las transiciones de solicitudes y reservas no se distribuyen en
condicionales de los controladores. \texttt{SolicitudReservaState} y
\texttt{ReservaState} definen los contratos; sus estados concretos encapsulan
las acciones permitidas, mientras \texttt{SolicitudReservaServiceImpl} y
\texttt{ReservaServiceImpl} delegan en ellos. El patrón centraliza invariantes y
facilita probar transiciones, a cambio de mantener más clases y actualizar el
selector cuando se incorpora un estado.

\subsection{Observer, Facade y Singleton}

\textbf{Observer.} La creación o el cambio de estado de un perfil debe poder
producir efectos secundarios sin acoplarlos al caso de uso.
\texttt{PerfilEventListener} define el observador,
\texttt{LoggingPerfilEventListener} implementa el registro y
\texttt{PerfilServiceImpl} publica \texttt{PerfilEvent}. La extensión con otros
observadores no modifica el servicio; como consecuencia crítica, la notificación
actual es síncrona y un listener costoso podría aumentar la latencia.

\textbf{Facade.} Construir el detalle completo de un laboratorio requiere
coordinar laboratorio, piso, bloque, campus y equipos.
\texttt{LaboratorioDetalleFacadeImpl}, a través de
\texttt{LaboratorioDetalleFacade}, ofrece una operación única y evita que el
controlador replique esa coordinación. El beneficio es una interfaz simple para
una composición real; la implementación mantiene, no obstante, una dependencia
con un DTO de presentación que constituye otra oportunidad de desacoplamiento.

\textbf{Singleton.} El conteo de peticiones HTTP requiere un registro único por
aplicación para no fragmentar las métricas. Las implementaciones
\texttt{HttpRequestsMetricsRegistry} de Académico y Usuarios son componentes de
Spring, conservan una instancia compartida y la exponen mediante
\texttt{getInstance()}; los filtros HTTP son sus clientes. Se obtiene un punto
de registro coherente, pero se mezcla el ciclo de vida del contenedor con estado
estático, lo que dificulta su sustitución frente a una inyección de constructor
convencional.

En conjunto, los siete patrones resuelven persistencia, creación, variación de
políticas, transiciones, reacción a eventos, composición y registro compartido.
Su valor no reside en maximizar el número de patrones, sino en hacer explícitas
decisiones que reducen acoplamiento en rutas verificables. Las limitaciones
descritas muestran que el diseño sigue siendo perfectible y proporcionan una
base concreta para trabajo posterior.
